\section{落地附录：把访谈观点转成可执行清单}

\subsection{附录 A：通用 Agent 产品发现检查表}

\begin{center}
\begin{tabular}{p{3.1cm}p{3.4cm}p{5.8cm}p{2.2cm}}
\toprule
检查项 & 关键问题 & 通过标准 & 优先级 \\
\midrule
任务定义 & 用户最终交付物是否明确 & 任务目标可在一句话内描述 & P0 \\
约束识别 & 权限/时限/合规是否明确 & 有明确边界且可程序化校验 & P0 \\
长尾价值 & 是否存在现有工具难解场景 & 至少 30\% 样本需跨工具编排 & P0 \\
可解释失败 & 失败是否能回放与定位 & 保留完整执行轨迹与日志 & P0 \\
交付质量 & 输出能否直接进入工作流 & 可执行、可复核、可复用 & P1 \\
成本结构 & 算力与人工成本是否可控 & 单任务成本随规模下降 & P1 \\
迭代节奏 & 是否建立周级复盘机制 & 每周有失败清单和修复清单 & P1 \\
组织对齐 & 目标是否被团队共享 & 所有任务挂到同一目标树 & P1 \\
\bottomrule
\end{tabular}
\end{center}

\subsection{附录 B：系统架构评审清单}

\begin{importantbox}{架构评审前必须回答的 10 个问题}
\begin{enumerate}
    \item 执行环境是否隔离，是否支持长任务持续运行
    \item 任务状态是否可观测、可回放、可审计
    \item 失败后是否存在自动恢复或安全中止机制
    \item 工具调用策略是否可解释，是否有降级路径
    \item 模型输出与工具结果是否有一致性检查
    \item 是否区分“回答正确”与“任务完成”
    \item 任务链路是否存在单点权限风险
    \item 是否支持灰度发布与回滚
    \item 新功能是否增加系统不可控复杂度
    \item 本次迭代是否明确了“不做什么”
\end{enumerate}
\end{importantbox}

\subsection{附录 C：评测日报模板（示例）}

\begin{center}
\begin{tabular}{p{2.8cm}p{2.5cm}p{2.5cm}p{2.8cm}p{3.5cm}}
\toprule
日期 & Completion Rate & Retry Cost & 人工介入率 & 主要失败类型 \\
\midrule
Day 1 & 61\% & 2.4 次/任务 & 38\% & Tool 误选、状态漂移 \\
Day 2 & 64\% & 2.2 次/任务 & 34\% & GUI 解析错误 \\
Day 3 & 66\% & 2.0 次/任务 & 31\% & 长任务超时 \\
Day 4 & 69\% & 1.8 次/任务 & 29\% & 输出格式不稳定 \\
Day 5 & 71\% & 1.7 次/任务 & 26\% & 边界权限处理 \\
\bottomrule
\end{tabular}
\end{center}

上表只是格式示例，不代表访谈中披露的真实内部数据。它的作用是强调：团队必须把“任务完成、重试成本、人工介入”作为同一面板上的联动指标，而不是孤立观察某一个分数。

\begin{warningbox}{只看单指标会导致错误优化}
例如只追求 completion rate，可能通过激进策略提高短期通过率，却显著增加重试成本和人工介入率，最终拖累真实用户体验。
\end{warningbox}

\subsection{附录 D：30 天迭代路线图（可直接复用）}

\begin{knowledgebox}{第 1 周：建立可观测性}
\begin{itemize}
    \item 打通任务日志、工具日志、错误日志
    \item 定义失败分类标签（推理、工具、执行、交付）
    \item 选择 20 个高价值真实任务做基线评测
\end{itemize}
\end{knowledgebox}

\begin{knowledgebox}{第 2 周：修补高频失败模式}
\begin{itemize}
    \item 针对 Top 3 失败类型做策略修补
    \item 为高频任务沉淀模板与自检步骤
    \item 建立灰度发布与回滚机制
\end{itemize}
\end{knowledgebox}

\begin{knowledgebox}{第 3 周：提高交付质量与用户感知}
\begin{itemize}
    \item 优化输出结构、可读性与复用性
    \item 引入人工审阅回路校正审美和叙事问题
    \item 对关键任务增加“完成后自检”流程
\end{itemize}
\end{knowledgebox}

\begin{knowledgebox}{第 4 周：固化飞轮并开始削减复杂度}
\begin{itemize}
    \item 复盘新增功能的真实收益
    \item 移除低复用工具与冗余流程
    \item 发布下一周期“不做清单”
\end{itemize}
\end{knowledgebox}

\subsection{附录 E：访谈金句与工程解释}

\begin{center}
\begin{tabular}{p{5.4cm}p{7.2cm}}
\toprule
访谈原意（转述） & 工程解释 \\
\midrule
“先做再说” & 先获取新信息增量，再做抽象优化 \\
“不做什么更重要” & 资源有限时，删减比扩张更能提升系统稳定性 \\
“长尾才有惊喜” & 头部任务易同质化，长尾任务更能拉开差距 \\
“Coding 是灵魂” & 统一动作空间，降低策略与维护复杂度 \\
“模型无法内化环境” & Agent 层长期存在，不会被模型能力完全替代 \\
\bottomrule
\end{tabular}
\end{center}

\subsection{本章小结}

附录的目标是把访谈观点变成可执行动作：先定义任务，再做架构；先做可观测，再做优化；先建飞轮，再扩能力。对任何通用 Agent 团队，这都是可落地且可复用的起点。


